\documentclass[10pt,a4paper]{amsart}
\usepackage[margin=19mm]{geometry}
\usepackage{amsmath,amssymb,mathtools}
\usepackage[scr=rsfs,cal=euler]{mathalfa}
\usepackage{xfrac}
\usepackage[unicode,hidelinks,bookmarks=false]{hyperref}
\newcommand{\ie}{i.e.\ }
\newcommand{\cf}{cf.\ }
\newcommand{\resp}{respectively\ }
\newcommand{\Subsection}{\subsection}
\providecommand{\nobreakdash}{\nobreak\hbox{-}}
\setlength{\emergencystretch}{2em}
\multlinegap0pt
\usepackage{amssymb}
\usepackage{mathtools}
\usepackage[shortlabels]{enumitem}
\usepackage{xcolor}
\usepackage{tikz}
\newtheorem{prop}{Proposition}
\newtheorem{lem}{Lemma}
\usepackage{cleveref}
\crefname{prop}{Proposition}{Propositions}
\crefname{lem}{Lemma}{Lemmas}
\newcommand{\N}{\mathbb{N}}
\newcommand{\dist}{\mathcal{D}}
\newcommand{\dx}{\mathrm{d}}
\newcommand{\finiteset}[1]{[#1]}	% finite set {1, .., n} or perhaps {0, .., n-1}, depending on which convention is more convenient
\newcommand{\two}{\finiteset{2}}		% two-element set
\title{Naturality characterizes product measures and relative product measures}
\author{Victor Bloch, Tobias Fritz}
\date{Source: arXiv:2609.15909v1 (CC BY 4.0)}
\begin{document}
\maketitle

\noindent\emph{Source and licence.} Naturality characterizes product measures and relative product measures. Version arXiv:2609.15909v1 (CC BY 4.0), distributed by arXiv under the Creative Commons Attribution 4.0 International licence, http://creativecommons.org/licenses/by/4.0/, as recorded in the arXiv licence metadata for this exact version and shown on the abstract page https://arxiv.org/abs/2609.15909v1. Full licence notice: \texttt{licenses/arxiv-2609.15909-CC-BY-4.0.txt}.\par\smallskip

\noindent\textit{Editorial scope.} Counted unit: the article's lem lem:relative\_bernoulliirrational (source0.tex lines 960--966). The article's complete original proof of it is delivered with it (source0.tex lines 968--1021, one \texttt{proof} environment of 54 lines). No mathematics is written, repaired or re-derived: the statement and the proof are the article's own lines, in the article's own notation, and no retained line is substituted or re-flowed.  Retained uncounted as the source-local context the proof needs: lines 879--891.  Nothing was omitted from any retained span, so the package carries no omission record.  No retained line contains a citation, so the wrapper carries no bibliography and no bibliography entry.  No retained line contains a figure environment, an image-inclusion command or drawing code, so no image is delivered and the package declares no asset dependency.  Editorial formatting only, in addition: an AMS \texttt{amsart} preamble that restates the article's own declarations and macros used by the retained lines, namely the environments \texttt{prop}, \texttt{lem} on separate counters, together with the standard packages those lines need.  The retained environments print in this document as Proposition 1, Lemma 1 in order of appearance, whereas the article prints the same items with the numbering of its own full document.  The complete native source of the article as distributed in the arXiv e-print for this exact version is bundled unchanged as \texttt{sources/210465/source0.tex}.
\begin{prop}\label{prop:relative_rationalcase}
	Let $X,Y$ be nonempty finite sets and $p : B \to \dist X$, $q : B \to \dist Y$ Markov kernels.
	Suppose that there are positive integers $N,M$ such that
	\[
		Np(x|b) \in \N, \qquad Mq(y|b) \in \N
		\qquad\quad
		\forall x \in X, y \in Y, b \in B.
	\]
	Then
	\[
		\alpha(\overline{p},\overline{q}) = \overline{p} \otimes_\omega \overline{q}.
	\]
\end{prop}

\begin{lem}\label{lem:relative_bernoulliirrational}
	For all Markov kernels $p,q : B \to \dist\two$, we have
	\[
		\alpha(\overline{p},\overline{q})
		= \overline{p} \otimes_\omega \overline{q}.
	\]
\end{lem}

\begin{proof}
	We first assume that $q$ has a common denominator as in \cref{prop:relative_rationalcase}.
	For $T \subseteq \two$ and $S \in \Sigma_B$, define
	\[
		\phi_{T,S}(s) \coloneq
		\alpha(\overline{s},\overline{q})(\{1\} \times T \times S).
	\]
	If $p \leq p'$, let $\eta$ be the witness above and put
	$r \coloneq \alpha(\overline{\eta},\overline{q})$.
	Naturality under $\theta_B^\pm$ gives
	\begin{align*}
		\phi_{T,S}(p')-\phi_{T,S}(p)
		&= r(\{1,2\} \times T \times S)-r(\{1\} \times T \times S) \\
		&= r(\{2\} \times T \times S) \geq 0.
	\end{align*}
	Thus $\phi_{T,S}$ is monotonically nondecreasing.

	For an arbitrary kernel $p$, define measurable dyadic approximations by
	\[
		p_n^-(1|b) \coloneq 2^{-n}\lfloor 2^n p(1|b)\rfloor,
		\qquad
		p_n^+(1|b) \coloneq 2^{-n}\lceil 2^n p(1|b)\rceil,
	\]
	and $p_n^\pm(2|b) \coloneq 1-p_n^\pm(1|b)$.
	These satisfy $p_n^- \nearrow p$ and $p_n^+ \searrow p$ pointwise, and both have common denominator $2^n$.
	By monotonicity and \cref{prop:relative_rationalcase},
	\[
		\int_S p_n^-(1|b) \, q(T|b) \, \omega(\dx b)
		\leq \phi_{T,S}(p)
		\leq \int_S p_n^+(1|b) \, q(T|b)\,\omega(\dx b).
	\]
	Taking $n \to \infty$ and using dominated convergence gives
	\[
		\phi_{T,S}(p) = \int_S p(1|b) \, q(T|b) \, \omega(\dx b).
	\]
	We likewise get the same identity with $\{2\}$ in place of $\{1\}$.
	Since evaluation on sets of the form $\{i\} \times \{j\} \times S$ distinguishes measures on $\two \times \two \times B$,
	the claim follows for arbitrary $p$ whenever $q$ has a common denominator.

	Now fix arbitrary $p$ and $q$.
	Applying the same witness construction in the second variable shows that, for every $T \subseteq \two$ and $S \in \Sigma_B$, the map
	\begin{align*}
		\dist\two & \longrightarrow [0,1] \\
		t & \longmapsto \alpha(\overline{p},\overline{t})(T \times \{1\} \times S)
	\end{align*}
	is monotonically nondecreasing.
	Approximate $q$ from below and above by $q_n^-$ and $q_n^+$ as above.
	The case already proved and dominated convergence give
	\[
		\alpha(\overline{p},\overline{q})(T \times \{1\} \times S)
		= \int_S p(T|b) \, q(1|b) \, \omega(\dx b).
	\]
	Again we get the same identity with $\{2\}$ in place of $\{1\}$ as well, proving the claim.
\end{proof}

\end{document}
